\documentclass[a4paper,11pt]{article}

\usepackage[ngerman]{babel}
\usepackage[top=3cm,bottom=3cm,left=3cm,right=3cm]{geometry} %NICHT IN KEMIE2


%\usepackage{microtype} %schöneres typesetting  mit [stretch=10] für nur 1% anstelle von 2%
%\usepackage{lmodern} %Latin Modern FONT
%\usepackage{pifont} %schönere Font



\usepackage{amsmath}
\usepackage{nicefrac}
\usepackage{amssymb}
%\usepackage{mathtools}

\usepackage{titlesec} %Um Section, etc bearbeiten zu können
\usepackage{setspace}
\usepackage{enumitem} %enumerate
\usepackage{multicol}
\usepackage[many]{tcolorbox}
\usepackage{hyperref}
\usepackage{arydshln}

%\renewcommand{\ttdefault}{lmtt} %Schriftart wird geändert zu Latin Modern Typewriter


\hypersetup{hidelinks,
			colorlinks=true,
			linkcolor=black,
			citecolor=miudiutex,
			urlcolor=miugray2
			}

\setlength{\parindent}{0pt}
\setcounter{section}{0}


\titlespacing*{\section}{0pt}{20mm}{6mm}
\titlespacing*{\subsection}{0pt}{6mm}{3mm}
\titlespacing*{\subsubsection}{0pt}{3mm}{0mm}



%\definecolor{miudiutex}{HTML}{2f3d39}
\definecolor{miudiutex}{HTML}{1d2926}
\definecolor{miugray}{HTML}{b4cccf}
\definecolor{miugray2}{HTML}{4d7365}
\definecolor{red}{HTML}{cf1f5c}
\definecolor{green}{HTML}{00A36C}
\definecolor{delphinidin}{HTML}{315794}
\definecolor{pelargonidin}{HTML}{b33807}
\definecolor{cyanidin}{HTML}{ba0746}
\definecolor{petunidin}{HTML}{b56fed}



\raggedcolumns %wichtig für Multicolumns


%================================
% SYMBOLS
%================================

\newcommand{\RA}{$\rightarrow$~}
\newcommand{\RL}{$\rightleftharpoons$~}
\newcommand{\HARP}{\ding{226}~}
%$\xrightarrow{2}$ %Pfeil mit 2 drüber


%================================
% SHORT COMMANDS (BOXES ; ETC)
%================================


\newcommand{\notiz}[1]{\begin{notizz}{}{}\begin{center}
			{\color{miugray2!50!miudiutex}#1} \end{center}\end{notizz}}
\newcommand{\zit}[2]{\begin{zitat*}{#1}{}\small #2 \end{zitat*}}

\newcommand{\frage}[2]{\begin{qst*}{#1}{}\small #2 \end{qst*}}

\newcommand{\unten}{\end{center} \tcblower \begin{center}}


\newcommand*\circled[1]{\tikz[baseline=(char.base)]{
            \node[shape=circle,draw,inner sep=0.5pt,miudiutex] (char) {\tiny #1};}}
            
\newcommand{\miusquare}{{\color{miugray2!50}\scriptsize $\blacksquare$}~}
\newcommand{\miusquarp}{{\color{petunidin!50}\scriptsize $\blacktriangleright$}~}
\newcommand{\niu}{\item[\miusquare]}
\newcommand{\nip}{\item[\miusquarp]}

%%%%% ABSAETZE

\newcommand{\pps}{\par \vspace*{1mm}}
\newcommand{\pp}{\par \vspace*{3mm}}
\newcommand{\ppbig}{\par \vspace*{8mm}}
\newcommand{\pphuge}{\par \vspace*{10mm}}

%================================
% FRAGE BOX
%================================

\tcbuselibrary{theorems,skins,hooks}
\newtcbtheorem{qst}{Frage:}
{%
	enhanced,
	breakable,
	colback = gray!15!white,
	frame hidden,
	boxrule = 0sp,
	borderline west = {2.5pt}{0pt}{red!70},
	sharp corners,
	detach title,
	before upper = \tcbtitle\par\smallskip,
	coltitle = black,
	fonttitle = \bfseries\sffamily,
	description font = \mdseries,
	separator sign none,
	segmentation style={solid, gray!50!black},
}
{th}




%================================
% ZITAT BOX
%================================

\tcbuselibrary{theorems,skins,hooks}
\newtcbtheorem{zitat}{Zitat}
{%
	enhanced,
	breakable,
	colback = gray!15!white,
	frame hidden,
	boxrule = 0sp,
	borderline west = {2.5pt}{0pt}{black!70},
	sharp corners,
	detach title,
%	before upper = \tcbtitle\par\smallskip,
	coltitle = gray!50!black,
	fonttitle = \bfseries\sffamily,
	description font = \mdseries,
%	separator sign none,
	segmentation style={solid, gray!50!black},
}
{th}



%================================
% SIMPLE SCHWARZER RAHMEN BOX
%================================


\newcommand{\boxfr}[1]{
	\begin{tcolorbox}[
	drop shadow southeast,
	enhanced,
	colframe=black!50,
	colback=white,
	boxrule = 1pt, 
%	toprule=0pt, bottomrule=0pt,rightrule=0pt,leftrule=0pt,
	boxsep=5pt,
	arc=1pt,
	]
	\begin{center}
	#1
	\end{center}
	\end{tcolorbox}
}

\definecolor{miudiutex}{HTML}{1d2926}
\definecolor{miugray}{HTML}{b4cccf}
\definecolor{miugray2}{HTML}{4d7365}
\definecolor{white2}{HTML}{f5f7f8}
\definecolor{red}{HTML}{cf1f5c}
\definecolor{green}{HTML}{00A36C}
\definecolor{delphinidin}{HTML}{315794}
\definecolor{uniblau}{HTML}{004e9f}
\definecolor{unigelb}{HTML}{fcba00}
\definecolor{unigrau}{HTML}{909085}
\definecolor{pelargonidin}{HTML}{b33807}
\definecolor{cyanidin}{HTML}{ba0746}
\definecolor{paeonidin}{HTML}{d15ca6}
\definecolor{petunidin}{HTML}{b56fed}
\definecolor{malvidin}{HTML}{8a2d8c}

\newcommand{\metermilli}{\mskip3mu \text{mm}}
\newcommand{\cm}{\mskip3mu \text{cm}}
\newcommand{\m}{ \text{m}}
\newcommand{\lite}{ \text{l}}
\newcommand{\kg}{ \text{kg}}
\newcommand{\g}{ \text{g}}
\newcommand{\km}{ \text{km}}
\newcommand{\h}{ \text{h}}
\newcommand{\s}{ \text{s}}
\newcommand{\tonne}{ \text{t}}
\usepackage[utf8]{inputenc}
\usepackage[T1]{fontenc}
\usepackage{lmodern}

\usepackage[
    activate={true,nocompatibility},
    final,
    tracking=true,
    kerning=true,
    spacing=true,
    factor=1100,
    stretch=30,
    shrink=20
]{microtype}
\usepackage{pdfpages}
\usepackage{awesomebox}
\usepackage{marginnote}
\tcbuselibrary{documentation}
\usepackage{sidenotes}
\usepackage{import}
\usepackage{xifthen}
\usepackage{transparent}
\usepackage[table]{xcolor}


\newcommand{\abbicon}{%
  \protect\tikz[baseline=0.2mm,scale=0.8]{%
    % Das blaue Quadrat mit abgerundeten Ecken
    \fill[uniblau, rounded corners=1.2pt] (0,0) rectangle (0.8em, 0.8em);
    % Der weiße Kreis in der Mitte
    \fill[white] (0.4em, 0.4em) circle (0.12em);
  }%
} 
\usepackage[labelfont={bf},font={color=miudiutex,small},figurename={\abbicon $~$Abbildung}]{caption}


\newcommand{\bckgrnd}{
\AddToHookNext{shipout/background}{
\begin{tikzpicture}[remember picture, overlay]
 \draw[draw=none,fill=uniblau!70, shift={(current page.north west)}] (-3,0) rectangle (23,-3.75);
 \draw[draw=none,fill=miugray, shift={(current page.north east)}] (-7,0) rectangle (3,-3.75);
\end{tikzpicture}
}
}



\newcommand{\incfig}[2]{%
\def\svgwidth{#2\columnwidth}
\import{C://LaTeX-Files/pngs/}{#1.pdf_tex}
}

\renewcommand{\textbf}[1]{\textsf{{\bfseries {#1}}}}



\titleformat{\section}[hang]
		{\LARGE \color{uniblau!75!black} \sffamily \bfseries}
		{\thesection}
		{6mm}
		{}
		[]
\titleformat{\subsection}[hang]
		{\large \sffamily \bfseries \color{black}}
		{\thesubsection}
		{4mm}
		{{\color{miugray}$\blacksquare ~$}}

\titleformat{\subsubsection}[block]
		{\normalsize \bfseries \sffamily \color{miudiutex}}%
		{\normalsize \color{black} \thesubsubsection}
		{2mm}
		{{\color{miugray}$\blacksquare ~$}}
		[ \color{black}]
		
\titleformat{\paragraph}
		{\normalsize \bfseries \sffamily \color{black}}
		{\footnotesize \theparagraph}
		{1mm}
		{}

\pagestyle{empty}
\titlespacing*{\section}{0pt}{20mm}{12mm}
\titlespacing{\section}{0pt}{20mm}{12mm}

\titlespacing*{\subsection}{0pt}{14mm}{4mm}
\titlespacing{\subsection}{0pt}{14mm}{4mm}

\titlespacing*{\subsubsection}{0pt}{6mm}{3mm}
\titlespacing{\subsubsection}{0pt}{6mm}{3mm}

\titlespacing*{\paragraph}{0pt}{6mm}{2mm}
\titlespacing{\paragraph}{0pt}{6mm}{2mm}


\let\oldmarginfigure\marginfigure
\let\endoldmarginfigure\endmarginfigure

\let\oldmarginfigure\marginfigure
\let\endoldmarginfigure\endmarginfigure

\renewenvironment{marginfigure}[1][0pt] % [1] steht für 1 Argument, [0pt] ist der Standardwert
  {\oldmarginfigure[#1]\captionsetup{labelfont={sf,bf,color=uniblau!75!black},font={color=miudiutex,small}}}
  {\endoldmarginfigure}
  
  
  
  
\renewcommand{\labelitemi}{\color{uniblau!70}$\bullet$}
\renewcommand{\labelitemii}{\color{uniblau!70}$\cdot$}
\newcommand{\plmtsquare}{{\color{uniblau!70}$\blacksquare~$}}
\newcommand{\splmt}[1]{{\color{uniblau!75!black} \sffamily #1}}
\newcommand{\fplmt}[1]{{\color{uniblau!75!black} \sffamily \bfseries #1}}

\newcommand{\antwort}[1]{
\begin{tcolorbox}[
	enhanced,
	enforce breakable,
	standard jigsaw,
	boxrule=0.7pt,
	colframe=black,
	sharp corners,
	boxsep=5pt,
	title=\bfseries \sffamily Antwort,
	opacityback=0,
	coltitle=miudiutex,
	attach boxed title to top text left={yshift=-\tcboxedtitleheight/2},
	boxed title style={colback=white,colframe=white},
	bottomrule at break=0pt,
	toprule at break=0pt,
	pad at break=16mm,
]
	\begin{center}
	#1
	\end{center}
	\end{tcolorbox}
}




\rowcolors{2}{miugray!50}{miugray!50}
\arrayrulecolor{white} % Weiße Gitterlinien
\setlength{\arrayrulewidth}{1.5pt} % Etwas dickere Linien
\renewcommand{\arraystretch}{1.4} % Mehr Zeilenhöhe für bessere Lesbarkeit}


%\setlist[itemize]{itemsep=0mm}
\setlist[enumerate]{itemsep=0mm}




\newif\ifshowme
%\showmetrue




\begin{document}
\newgeometry{top=1.8cm,bottom=4cm,left=1.5cm,right=7.5cm,marginparsep=-2.00cm,marginparwidth=4.75cm}


\bckgrnd


\begin{center}
\LARGE
\color{white}
\textbf{Übung 12} \par \vspace*{2mm}
\large \color{miugray} \textbf{Prozessbezogene Lebensmitteltechnologie}
\end{center}
\par 
\vspace*{6mm}


\color{miudiutex}
\section*{Sinkgeschwindigkeit}

\subsection*{Aufgabe 1}
Zwischen dem Serum und den Fettkügelchen von Rohmilch besteht ein Dichteunterschied
von \mbox{$\Delta \rho$ = 99,8 kg/m$^3$.} Die Viskosität der Milch beträgt $\eta$ = 1,79 $\cdot$ 10$^{-3}$ Pa $\cdot$ s. Die Dichte ist mit \mbox{$\rho$ = 1.028 kg/m$^3$} angegeben. 

\begin{enumerate}[label=\alph*)]

\item Berechnen Sie die Aufrahmgeschwindigkeit eines Fettkügelchens mit einem
Durchmesser von \mbox{4 µm} unter der Annahme einer laminaren Strömung.

\ifshowme
\antwort{
\begin{align*}
v &= \frac{d^2 \cdot g \cdot (\rho_1 - \rho_2)}{18 \cdot \eta} \\
v &= \frac{(4 \cdot 10^{-6} ~\text{m})^2 \cdot 9,81 ~\text{m/s}^2 \cdot 99,8 ~\text{kg/m}^3}{18 \cdot 1,79 \cdot 10^{-3} ~\text{Pa s}} \\
v &= 4,86 \cdot 10^{-7} ~\text{m/s}
\end{align*}
}
\else
\fi




\item Prüfen Sie mit Hilfe der Reynoldszahl, ob die Annahme einer laminaren Strömung
korrekt war. Erinnerung: Eine laminare Strömung kann angenommen werden, wenn
die Reynoldszahl $<$ 0,5 ist.
\ifshowme
\antwort{
\begin{align*}
Re &= \frac{\rho \cdot d \cdot v}{\eta} \\
Re &= \frac{1028 ~\text{kg/m}^3 \cdot 4 \cdot 10^{-6} ~\text{m} \cdot 4,86 \cdot 10^{-7} ~\text{m/s}}{1,79 \cdot 10^{-3} ~\text{Pa s}} \\
Re &= 1,116 \cdot 10^{-6}
\end{align*}
\RA Die Annahme einer laminaren Strömung war korrekt

}
\else
\fi


\item Die Fettkügelchen (d = 4 µm) haben nach 7.500 Minuten die Oberfläche erreicht. In
welcher Zeit schaffen es die größten Fettkügelchen (d = 5 µm) an die Oberfläche?

\ifshowme
\antwort{

Strecke berechnen:
\begin{align*}
s = v \cdot t \rightarrow s &= 4,86 \cdot 10^{-7} ~\text{m/s} \cdot 7500 \cdot 60s \\
 &= 0,2187 ~\text{m} \\
\end{align*}
Geschwindigkeit größerer Fettkügelchen berechnen: 
\begin{align*}
v &= \frac{d^2 \cdot g \cdot (\rho_1 - \rho_2)}{18 \cdot \eta} \\
v &= \frac{(5 \cdot 10^{-6} ~\text{m})^2 \cdot 9,81 ~\text{m/s}^2 \cdot 99,8 ~\text{kg/m}^3}{18 \cdot 1,79 \cdot 10^{-3} ~\text{Pa s}} \\
v &= 7,59 \cdot 10^{-7} ~\text{m/s}
\end{align*}
Zeit die größere Fettkügelchen für Strecke $s$ benötigen:
\begin{align*}
t = \frac{s}{v} \rightarrow t &= \frac{0,2187 ~\text{m}}{7,59 \cdot 10^{-7} ~\text{m/s} \cdot 60} \\
&= 4802,4 ~\text{min}
\end{align*}
}
\else
\fi


\end{enumerate}

\subsection*{Aufgabe 2}
Aus 70 °C heißer Bierwürze sollen Heißtrub und Hopfentrester in einem runden Absetzbecken
(Füllhöhe 850 mm; Durchmesser 7000 mm) zur Vorklärung absedimentiert werden, bevor das
Gemisch zentrifugiert wird. Die Partikel sind zwischen 30 und 90 µm groß und haben eine
Dichte von 1.050 kg/m$^3$. Für die Dichte und die Viskosität der Würze bei 70 °C nehmen Sie bitte die Kennzahlen für Wasser und laminare Strömungsverhältnisse an.
\pp
Welche Partikel sind nach 40 min vollständig absedimentiert?


\ifshowme
\antwort{

Geschwindigkeit die Partikel brauchen um nach 40min vollständig absedimentiert zu sein:
\begin{align*}
v &= 0,85~\text{m} / ~2400~\text{s} \\
&= 3,5 \cdot 10^{-4} ~\text{m/s}
\end{align*}
Stokes umstellen:
\begin{align*}
v &= \frac{d^2 \cdot g \cdot (\rho_1 - \rho_2)}{18 \cdot \eta} \\
d &= \sqrt{\frac{v \cdot 18 \cdot \eta}{g \cdot \Delta p}} \\
d &= \sqrt{\frac{3,5 \cdot 10^{-4} ~\text{m/s} \cdot 18 \cdot 1 \cdot 10^{-3} ~\text{Pa s}}{9,81 ~\text{m/s} \cdot 50 ~\text{kg/m}^3}} \\
 &= 7,2 \cdot 10^{-5} ~\text{m} \\
 &= 72 ~\text{µm}
\end{align*}
Alle Fettkügelchen die größer als 72µm sind, sedimentieren vollständig.
}
\else
\fi


\pagebreak
\subsection*{Aufgabe 3}
Milch soll mittels Mikrofiltration entkeimt werden. Hierfür wählen Sie ein geeignetes
Filterelement mit einem Porendurchmesser von 0,2 µm, einer Länge von 20 cm und einer
Porenanzahl von 1.000.000. Gehen Sie dabei bitte von einer laminaren Strömung aus. Die
Viskosität der Milch beträgt 2 mPa$\cdot$s. Durch Anlegen eines Vakuums wird eine Druckdifferenz
von 0,5 bar erzeugt. Wie groß ist die Volumenstromstärke der Milch durch das Filterelement? 

\ifshowme
\antwort{
\begin{align*}
\dot{V} &= \frac{\pi \cdot \Delta p \cdot R^4}{8 \cdot l \cdot \eta} \\
\dot{V} &= \frac{\pi \cdot 5 \cdot 10^4 ~\text{Pa} \cdot (1 \cdot 10^{-7} ~\text{m})^4}{8 \cdot 0,2 ~\text{m} \cdot 2 \cdot 10^{-3} ~\text{Pa s}} \\
&= 4,91 \cdot 10^{-21} ~\text{m}^3/\text{s}
\end{align*}
Dadurch, dass wir aber eine Anzahl von 1.000.000 Poren haben, müssen wir den Volumenstrom $\dot{V}$ nochmal mit der Anzahl multiplizieren:
\begin{align*}
4,91 \cdot 10^{-21} ~\text{m}^3/\text{s} \cdot 1.000.000 = 4,91 \cdot 10^{-15}
\end{align*}

}
\else
\fi




\end{document}
